\subsubsection{Deutsch's Algorithm}\label{sec:deutschs-algorithm}

\paragraph{Problem Definition and the Four Possible Functions}\label{sec:problem-definition-and-the-four-possible-functions}

Before learning the algorithm itself, we first need to understand \textbf{what problem Deutsch's algorithm is trying to solve}.

\highspace
\begin{flushleft}
    \textcolor{Red2}{\faIcon{exclamation-triangle} \textbf{The Computational Problem}}
\end{flushleft}
Suppose someone gives us a mysterious device. We can feed it a bit ($0$ or $1$), and it returns another bit. We \textbf{cannot open the device} to see how it works, but we can ask questions by feeding it inputs and observing the outputs. This is what we call a \definition{Black Box} or \definition{Oracle}.

\highspace
Let's concretize this idea. Imagine the device belongs to another company. Maybe it is an encrypted circuit, a proprietary software routine, a remote server, or simply an expensive experiment. Every time we ask for an output, we pay a cost. The cost could be computation time, money, energy, communication delay, or simply the number of allowed queries (for example, if the device is a quantum experiment that can only be run a limited number of times). Therefore, the \textbf{objective is \underline{not} to reconstruct the function}. Instead, we want to learn \textbf{some useful property of the function while asking as few questions as possible}. This is the essence of a \definition{Query Complexity Problem}.

\highspace
Suppose we are given a function:
\begin{equation*}
    f : \{0, 1\} \to \{0, 1\}
\end{equation*}
This notation simply means:
\begin{itemize}
    \item The input is \textbf{one classical bit};
    \item The output is \textbf{one classical bit}.
\end{itemize}
Unlike previous examples, \hl{we are not interested in computing the value of the function.} Instead, we only want to discover one global property of the function.

\highspace
Since there are only two possible inputs, there are only two function evaluations:
\begin{equation*}
    f(0) \in \{0, 1\} \quad \text{and} \quad f(1) \in \{0, 1\}
\end{equation*}
Deutsch's Algorithm \textbf{classifies the function} into one of \textbf{two categories}:
\begin{itemize}
    \item \definition{Constant Function}. A function is \emph{constant} if it \textbf{always returns the same output}. Mathematically:
    \begin{equation*}
        f(0) = f(1)
    \end{equation*}
    There are exactly two constant functions:
    \begin{itemize}
        \item $f(x) = 0$
        \item $f(x) = 1$
    \end{itemize}
    In other words, the function always returns 0 or always returns 1, regardless of the input. Their truth tables are:
    \begin{center}
        \begin{tabular}{@{} c c c @{}}
            \toprule
            $x$ & $f(x)=0$ & $f(x)=1$ \\
            \midrule
            0 & 0 & 1 \\
            1 & 0 & 1 \\
            \bottomrule
        \end{tabular}
    \end{center}
    Notice that the output never changes, regardless of the input.


    \item \definition{Balanced Function}. A function is \emph{balanced} if it returns:
    \begin{itemize}
        \item \textbf{0 exactly once}, and
        \item \textbf{1 exactly once}.
    \end{itemize}
    Since there are only two possible inputs, this means:
    \begin{equation*}
        f(0) \neq f(1)
    \end{equation*}
    Again, there are exactly two balanced functions:
    \begin{itemize}
        \item $f(x) = x$
        \item $f(x) = \overline{x}$
    \end{itemize}
    Where $\overline{x}$ is the logical NOT of $x$. Their truth tables are:
    \begin{center}
        \begin{tabular}{@{} c c c @{}}
            \toprule
            $x$ & $f(x)=x$ & $f(x)=\overline{x}$ \\
            \midrule
            0 & 0 & 1 \\
            1 & 1 & 0 \\
            \bottomrule
        \end{tabular}
    \end{center}
    Unlike constant functions, balanced functions return each output exactly once.
\end{itemize}
So, instead of identifying exactly which one of the four functions we have, Deutsch proposed an easier question: \emph{\textbf{is the function constant or balanced?}} That sounds almost trivial, but notice something important. To answer this question, we don't need every detail about the function. We only want one \textbf{global property} (constant or balanced) of the function.

\highspace
In other terms, Deutsch's algorithm wants to answer a deeper question: \emph{\textbf{can a quantum computer determine a global property of an unknown function using fewer queries than any classical computer?}}

\highspace
\begin{definitionbox}[: Deutsch's Algorithm]
    \definition{Deutsch's Algorithm} is the first quantum algorithm to demonstrate a genuine quantum computational advantage. Given \textbf{oracle} access to an unknown Boolean function (black box):
    \begin{equation*}
        f : \{0, 1\} \to \{0, 1\}
    \end{equation*}
    Promised to be either \textbf{constant} or \textbf{balanced}, the algorithm determines which of the two cases holds using \textbf{a single query} to the quantum oracle. A deterministic classical algorithm requires \textbf{two queries} to solve the same problem, because it must evaluate both $f(0)$ and $f(1)$ to determine whether the function is constant or balanced (we will analyze this in detail in the next section).
\end{definitionbox}

\highspace
\textcolor{Green3}{\faIcon{question-circle} \textbf{Why should we care today?}} The problem itself is not practically useful. Nobody builds quantum computers to classify one-bit functions. The importance of Deutsch's algorithm is that it introduces an idea that appears throughout quantum computing: \hl{rather than learning every output of a function, a quantum algorithm can exploit superposition, interference, and phase kickback to extract a global property directly.} This idea is reused by much more powerful algorithms like Simon's algorithm, Shor's algorithm, and Grover's algorithm. In fact, Deutsch's algorithm is often called the \textbf{\emph{Hello World} of quantum algorithms} because it is the simplest example of a quantum algorithm that outperforms classical algorithms.